\section{Separation of all the scalar sections on a signed mixed arc}

% BEGIN S83M1
\paragraph{S83M1. Inverse coordinates, the shift operator and the scalar residues.}
Use U83M1--U83M3, with actual cut vector $w$ at $c$. Write
$w=D_b a/\Lambda$, $a\in A^r$, in U83M6's translated basis.
Its two defining properties are $\Lambda w\in A\varphi^e L_{c-e}$
and $(\gamma-1)w\in vL_c$. There are integral constant numerator
coordinates $f=C_- a/\Lambda$ representing $p^\sigma w$.
Indeed apply U83M6's construction to $U_b^{-1}$ in place of
$U_b$. Its coefficient bounds are the same by J82F1. In the
quotient modulo $v\Lambda^2$, with $p$ inverted, $C_-$ is
$p^\sigma\varphi^e(U_b^{-1}\bmod v^M)$ and the two truncated
matrices are inverse before their scalar factors. Multiplication
by the forward truncated matrix therefore gives the required
$\mathrm{RepVV}(p^\sigma w,f)$ after clearing its coefficient
denominators. At each primitive level the Laurent expansions of
$d_bf$ and $p^\sigma w$ differ in degrees at least $H_s$ by the
inverse Taylor calculation of J82F2 with U83M6's modulus.
In particular $(\gamma-1)f\in v\mathcal N$. Indeed its numerator
after applying $\gamma-1$ vanishes at $v=0$. At every primitive
level division by $v$ is a Taylor unit and the two errors have
order at least $H_s$, so their difference has the same true
filtration test as $p^\sigma(\gamma-1)w/v\in L_c$.

Use $P=P_J$ of U83M4 and write $g=P(\gamma)f$. The integral
numerator condition survives since the factors $(P_s)$ are stable.
Moreover $g$ represents $p^\sigma P(\gamma)w$. For this last
identity put $B(v)=U_b\bmod v^M$. The gamma transform of $bB-bU_b$
agrees with zero through degree $M-1$, so $\gamma(bB)-bB$ has
coordinates in $v^M A[1/p]^r$ in the $b$-basis. After $e$
Frobenius steps the error is divisible by $f'$ and hence by
$v\Lambda^2$. Thus the rational congruence defining representation
is gamma-equivariant, also for denominator $\Lambda$ since its
ideal is stable. It is linear for polynomial combinations.

We verify the native single-coordinate conditions directly. At
level $s$ denote by $f_i$ and $g_i$ the coefficients in Laurent
degree $i$ using $1+X=\zeta\exp(Y)$. Write $F_s^a$ for the
true filtration. The coefficients of $(\gamma-1)f$ in degree
$i$ lie in $B_s\otimes F_s^{-i}$, because multiplication by
$v$ preserves the filtered Taylor lattice. For $i\ne0$ the
operator on the degree-$i$ quotient is $u^i\rho-1$, $u=1+p$,
where the branch permutation $\rho$ has order dividing $p^s$.
It is invertible, as seen by multiplying by the finite geometric
sum and using $u^{ip^s}\ne1$. Thus $f_i\in B_s\otimes F_s^{-i}$
for every nonzero degree. The same holds for $g_i$; and $g_i=0$
in each degree $i\in J$, by the factor $\gamma^{p^e}-u^{ip^e}$
in $P$. All these statements descend from a splitting field.

Take any one of the $U'$-coordinates of U83M3's splitting and call
its scalar $g_k$. Projection onto its basis vector along all other
coordinates kills $(S_r)_L$ and has image in every up subspace
$F_s^{b_s}$, $b_s=s_{c-s}$. At an up level all degrees $i<-b_s-S$
vanish, and the shift polynomial kills $-b_s-S\le i<-b_s$.
At a down level every $F_s^{-i}$ for $i<b_s$ is contained in
$F_s^{-\eta}\subset(S_r)_L$: for $-i<-\eta$ this follows from
the gap before the deep negative cloud. Hence every nonzero degree
$i<b_s$ projects to zero. The remaining $b_s\le i<b_s+S$ are
killed by $P$. At degree zero the class of $f_0$ modulo $F_s^0$
is fixed by both $\rho$ and the tame actions, using the same
filtration test for $(\gamma-1)f$ and tame invariance. Since
$F_s^0\subset(S_r)_L$ on this down letter, the chosen scalar
coordinate at zero is one value in $L$, common to every primitive
branch by K82P1. The same is true for $g_k$, since the polynomial
acts by its value at one on this projected constant.

Put $Q_+=\prod_{s\ \mathrm{up}}P_s^{b_s}$. These calculations give
$Q_+g_k\in A$: the numerator started integral with denominator
$\Lambda$, and all the unwanted primitive factors divide it with
the required multiplicities by distinguished division after
checking all roots. At any down level $s$ call its constant value
$z_s\in L$. Choose $\nu\ge0$ with $p^\nu z_s\in O$. The integral
numerator $Q_+(p^\nu g_k-p^\nu z_s)$ vanishes to order
$b_s+S$ at every primitive branch of that level, and hence is in
$P_s^{b_s+S} A$. This proves the entire native $\mathrm{InFF}$
for $g_k$, including the constant allowed to depend on $s$.
The basis vector itself is in the required up filtrations, exactly
$\mathrm{InLpp}$.
% END S83M1

% BEGIN S83M2
\paragraph{S83M2. Reduction of the mixed row to the individual scalar lifts.}
Keep the fixed cut and let $g_k e_k$ denote the vector of one
complementary scalar from S83M1 in constant coordinates. Then
$(\gamma-1)(g_k e_k)\in v\mathcal N$. At zero the numerator
of the difference is divisible by $v$. At up levels its Taylor
order is at least $-b_s$, and every coefficient is a multiple of
a vector in $F_s^{b_s}$. At down levels the scalar is constant
through degree $b_s+S-1$, with the same degree-zero value at
each branch, so the difference starts in degree $b_s+S$. At
those degrees $F_s^{-i}=D_c$. Division by the Taylor unit $v$
preserves the conditions throughout. Also $(\gamma-1)g\in
v\mathcal N$, by S83M1 and stability of $v\mathcal N$ under
$\gamma$. Subtract all complementary coordinates. The resulting
$g_r$ is supported on $S_r$ and belongs to the fixed quotient
of U83M3.

Let $E_1=\lfloor40U+\eta I\rfloor$; this is at least $\sigma$
by U83M1 and $U\ge1$. Suppose for each complementary scalar
$g_k$ the native $\mathrm{LiftsAt}$ holds at $E_1$. It suffices
to know this only for the scalars arising in S83M1, for every
tested map and class. By definition it supplies an almost
invariant finite-pole source vector $w_{1k}$ and an element
$z_k\in vL_c(T_2)$ for which $F w_{1k}-z_k$ represents
$p^{E_1}g_k e_k$. This target difference is in the physical
space $\Lambda^{-1} A\varphi^e L_{c-e}(T_2)$: the native lift
places the first term there and $vL_c$ is contained there by
U83M2. U83M6 supplies a physical vector $V(g_r)$ representing
$p^\sigma g_r$. By linearity over $A$ the vector
\[
 y=p^{\sigma+E_1}P(\gamma)w-
        \sum_k(Fw_{1k}-z_k)-p^{E_1-\sigma}V(g_r)
\]
represents zero. For precision, all representation congruences
can be compared after inverting $p$ using one truncated matrix.
For zero numerator coordinates the assertion is
$\Lambda y\in v\Lambda^2 A\varphi^e L_{c-e}(T_2)$ after
inverting $p$. The left side is already integral. Distinguished
division on the free translated basis cancels the extra $p$-power,
giving $y\in v\Lambda A\varphi^e L_{c-e}(T_2)\subset vL_c(T_2)$.

Take the finite class attached to each physical vector as in
U83M6. The source vectors have finite poles and their generator
difference is in $vL_c(T_1)$ by $\mathrm{LiftsAt}$, so the
identical truncation and Frobenius contraction give natural
finite classes there too. A congruence of these vectors modulo
$vL_c$ gives equality of their classes by the same positive-grade
argument. Finally $P(\gamma)w-P(1)w\in vL_c$, since
$(\gamma-1)w\in vL_c$ and this module is stable. The scalar
$P(1)$ is $p^{\kappa_c}$ times a unit in $O$. Hence the cut
class at exponent $\sigma+E_1+\kappa_c$ lifts up to the
submodule generated by the actual exceptional classes of U83M6.
There are at most $me(e+1)$ of these, independent of the input
class. This verifies the entire cut comparison required in
U83M4 conditional only on those complementary single-coordinate
lifts. In particular there is no residual exceptional-module,
filtration-projection or simultaneous gluing hypothesis in the
mixed C/D row reduction.
% END S83M2
